% !TEX root = ../main.tex
%% Devoir surveillé n° 2 — chapitres 4 à 6
%% Résultats vérifiés par verifications/eval_ds.py.

\sujetEval{DS2}{MPSI · Devoir surveillé}{Devoir surveillé \numero2}
  {Durée : \textbf{4 heures}\\Calculatrice interdite\\Chapitres 4 à 6}{4}
  {Exercice 1/Fonctions usuelles/6 pts,Exercice 2/Équations différentielles/5 pts,%
   Exercice 3/Suite récurrente/5 pts,Exercice 4/{Continuité, dérivabilité}/4 pts}
\consigne{Les quatre exercices sont indépendants. On veillera à justifier
soigneusement les dérivabilités et les passages à la limite.}

\begin{exoEval}{6 points}{Fonctions usuelles}
\begin{questions}
  \item Montrer que pour tout $x\geqslant0$ :
    $x-\dfrac{x^{2}}{2}\leqslant\ln(1+x)\leqslant x$.
  \item Soit $\varphi(x)=\arctan x+\arctan\dfrac1x$ sur $\R^{*}$. Calculer
    $\varphi'$ et en déduire une expression simple de $\varphi$ sur
    $\intOO{0}{+\infty}$ et sur $\intOO{-\infty}{0}$.
  \item Étudier les variations de $h:x\mapsto x^{1/x}$ sur
    $\intOO{0}{+\infty}$. En déduire, sans calculatrice, lequel des deux
    nombres $\mathrm{e}^{\pi}$ et $\pi^{\mathrm{e}}$ est le plus grand.
\end{questions}
\end{exoEval}

\begin{exoEval}{5 points}{Équations différentielles}
\begin{questions}
  \item Résoudre sur $\R$ l'équation $(1+x^{2})\,y'+2x\,y=1$, puis
    déterminer la solution telle que $y(0)=1$. \emph{On pourra
    reconnaître une dérivée.}
  \item Résoudre sur $\R$ l'équation $y''-3y'+2y=\mathrm{e}^{x}$, puis
    déterminer la solution vérifiant $y(0)=y'(0)=0$.
\end{questions}
\end{exoEval}

\begin{exoEval}{5 points}{Suite récurrente}
On définit la suite $(u_n)$ par $u_0=0$ et $u_{n+1}=\sqrt{2+u_n}$.
\begin{questions}
  \item Montrer que $0\leqslant u_n<2$ et que $(u_n)$ est croissante. En
    déduire qu'elle converge, et déterminer sa limite.
  \item Montrer que $|u_{n+1}-2|\leqslant\frac12|u_n-2|$, puis que
    $|u_n-2|\leqslant\dfrac{2}{2^{n}}$.
  \item Montrer par récurrence que $u_n=2\cos\dfrac{\pi}{2^{n+1}}$.
  \item En déduire $\displaystyle\lim_{n\to+\infty}4^{n}(2-u_n)$.
\end{questions}
\end{exoEval}

\begin{exoEval}{4 points}{Continuité, dérivabilité}
\begin{questions}
  \item Soit $f$ continue sur $\intFF{0}{1}$ avec $f(0)=f(1)$. Montrer
    qu'il existe $c\in\intFF{0}{\frac12}$ tel que
    $f\bigl(c+\frac12\bigr)=f(c)$. \emph{Introduire
    $g(x)=f\bigl(x+\frac12\bigr)-f(x)$.} Interpréter géométriquement.
  \item Soit $P\in\R[X]$ possédant $p\geqslant2$ racines réelles distinctes.
    Montrer que $P'$ possède au moins $p-1$ racines réelles distinctes. En
    déduire que, pour $n\geqslant2$ et $a,b\in\R$, le polynôme
    $X^{n}+aX+b$ a au plus trois racines réelles distinctes.
\end{questions}
\end{exoEval}
